\section{Intuitionistic propositional logic by finite saturation}
\label{sec:ipc-saturation}
\providecommand{\IPC}{\mathsf{IPC}}
\providecommand{\ipra}{\mathsf{iPRA}}
\providecommand{\Forces}{\Vdash}

This section gives a direct, bounded proof of intuitionistic propositional
soundness, completeness, decidability, and the finite model property.
The construction uses neither an infinite maximal theory nor a preliminary
cut-elimination theorem. Its purpose is a transparent primitive recursive
formalization, rather than an efficient implementation. Kripke's original
semantics provides the historical setting \cite{kripke1965}; all bounds and
the finite construction used below are established here.

\subsection{The proof system and its finite restriction}

Formulas are generated from atoms $p_i$, falsity $\bot$, and the binary
connectives $\wedge,\vee,\to$. Write $\neg A$ for $A\to\bot$ and $\top$ for
$\bot\to\bot$. Contexts are finite sets. A judgement $\Gamma\vdash A$ has
one formula on the right. We use the following natural-deduction rules, with
weakening and cut displayed explicitly:
\[
\frac{A\in\Gamma}{\Gamma\vdash A}\quad
\frac{\Gamma\vdash A\quad\Gamma\subseteq\Delta}{\Delta\vdash A}\quad
\frac{\Gamma\vdash A\quad\Gamma\cup\{A\}\vdash B}{\Gamma\vdash B}\quad
\frac{\Gamma\vdash\bot}{\Gamma\vdash A}.
\]
\[
\frac{\Gamma\vdash A\quad\Gamma\vdash B}{\Gamma\vdash A\wedge B}
\qquad
\frac{\Gamma\vdash A\wedge B}{\Gamma\vdash A}
\qquad
\frac{\Gamma\vdash A\wedge B}{\Gamma\vdash B},
\]
\[
\frac{\Gamma\vdash A}{\Gamma\vdash A\vee B}\qquad
\frac{\Gamma\vdash B}{\Gamma\vdash A\vee B}\qquad
\frac{\Gamma\vdash A\vee B\quad
       \Gamma\cup\{A\}\vdash C\quad
       \Gamma\cup\{B\}\vdash C}{\Gamma\vdash C},
\]
\[
\frac{\Gamma\cup\{A\}\vdash B}{\Gamma\vdash A\to B}
\qquad
\frac{\Gamma\vdash A\to B\quad\Gamma\vdash A}{\Gamma\vdash B}.
\]
These are twelve rule forms, counting the two conjunction eliminations and
the two disjunction introductions separately. The displayed cut is merely
substitution of a derivation for an assumption, so this presentation proves
exactly ordinary intuitionistic natural deduction. More explicitly, a
derivation of $B$ from $\Gamma\cup\{A\}$ can be traversed recursively while
replacing each free use of $A$ by the supplied derivation of $A$; under a
discharged assumption, retain that local assumption and weaken the supplied
derivations. The same traversal handles the two discharged branches of
disjunction elimination. This is a primitive recursive transformation of
finite proof trees. It need not normalize the resulting proof.

Fix a finite input sequent $\Gamma_0\vdash A_0$. Let $S$ be the set of all
its subformulas, together with $\bot$, and enumerate its distinct members
as $s_0,\ldots,s_{m-1}$. Thus $m\geq1$, $\Gamma_0\subseteq S$, $A_0\in S$,
and $S$ is closed under taking immediate subformulas. A context is now an
$m$-bit mask. There are precisely
\[
                 N=m2^m
\]
possible judgements $(\Gamma,A)$ with $\Gamma\subseteq S$ and $A\in S$.
Restrict every displayed rule to formulas in $S$. The restriction is imposed
on all formulas, including cut formulas and the conclusions of introduction
rules.

Let $D_0$ contain precisely the assumption judgements. Having defined the
Boolean table $D_r$, let $D_{r+1}$ contain $D_r$ and every conclusion of a
restricted rule whose premises occur in $D_r$. All rule instances can be
enumerated by a bounded loop. Retain, for every newly added judgement, its
first stage and the first successful rule instance in a fixed ordering.

\begin{lemma}[Finite closure and proof certificates]
\label{lem:ipc-finite-closure}
The table $D=D_N$ is closed under all twelve restricted rules. Every entry
of $D$ has an $\IPC$ proof using only formulas in $S$. A simultaneous proof
certificate for all accepted judgements can be represented by a directed
acyclic graph with at most $N$ inference records. Expanding sharing gives
each accepted judgement a proof tree with fewer than $3^{N+1}$ nodes, with
cut retained as one of the displayed rules.
\end{lemma}
\begin{proof}
Every table is contained in its successor. If one transition adds no entry,
the defining update gives the same table forever. If $D_N$ were not closed,
each of the first $N+1$ transitions would therefore add at least one new
entry, which is impossible for a table with $N$ positions. This counting
argument is finite: count the set bits, and observe that a proper increase
increases the count by at least one.

For proof extraction, order accepted judgements by first appearance stage,
then by their numeric indices. An assumption entry is a leaf. Every other
record points to at most three records of strictly earlier stage and stores
the rule used. Thus the certificate is acyclic and has at most $N$ records.
Following the records backwards from an entry first appearing at stage $r$
produces a tree of height at most $r$, with at most three children per node.
Its node count is at most $1+3+\cdots+3^r<3^{r+1}$. Each node is one of the
displayed $\IPC$ rules, so this is an actual proof, not a certificate of
semantic plausibility.
\end{proof}

This also supplies explicit primitive recursive resource bounds. There are
at most $12N^4$ candidate tuples consisting of a rule tag, at most three
premise judgements, and a conclusion; fewer-premise rules may use dummy
entries. Hence $N$ rounds require at most $12N^5$ candidate checks. This is
a bound on checks, not a machine-step complexity claim about a particular
encoding. Each check is primitive recursive. An inference record uses a
four-bit rule tag and four indices bounded by $N$, so the graph can use
at most
\[
          N\bigl(4+4\lceil\log_2(N+1)\rceil\bigr)
\]
bits for those fields, in addition to the input list $S$ and elementary
length and root fields. A local checker recomputes the side conditions from
the referenced judgements. If an ordinary natural-deduction tree without
the explicit cut rule is desired, the substitution procedure above removes
the cuts. The stated $3^{N+1}$ bound is for expansion of sharing, not for
that additional transformation.

\subsection{Finite maximal extensions without a choice principle}

For $X\subseteq S$ and $C\in S$, abbreviate the corresponding table entry
by $D(X,C)$. Say that $X$ \emph{avoids} $C$ when $D(X,C)=0$. This is a
decidable property of two finite codes.

\begin{lemma}[Effective finite prime extension]
\label{lem:ipc-prime-extension}
If $X$ avoids $C$, a primitive recursive procedure produces $T\supseteq X$
such that:
\begin{enumerate}
\item $D(T,C)=0$;
\item for every $E\in S\setminus T$, $D(T\cup\{E\},C)=1$;
\item $D(T,E)=1$ implies $E\in T$ for every $E\in S$;
\item if $E\vee F\in T$, then $E\in T$ or $F\in T$;
\item $\bot\notin T$.
\end{enumerate}
\end{lemma}
\begin{proof}
Start with $T_0=X$. At stage $j<m$, test the bit
$D(T_j\cup\{s_j\},C)$. If it is zero, set
$T_{j+1}=T_j\cup\{s_j\}$; otherwise set $T_{j+1}=T_j$. Put $T=T_m$.
By construction, every $T_j$ avoids $C$, proving (1). If $E=s_j$ is absent
from $T$, it was rejected at stage $j$. The judgement
$T_j\cup\{E\}\vdash C$ belongs to $D$, and weakening gives
$D(T\cup\{E\},C)=1$. This proves (2).

For (3), decide whether $E\in T$. If it is absent but $D(T,E)=1$, combine
this judgement with (2) by the displayed cut rule. The result is
$D(T,C)=1$, contrary to (1). Thus $E\in T$. Notice that membership in $T$
is a bit test; this reasoning does not use excluded middle for an arbitrary
arithmetical proposition.

For (4), decide membership of $E$ and $F$. If neither belongs to $T$, (2)
gives both $D(T\cup\{E\},C)=1$ and $D(T\cup\{F\},C)=1$. The assumption
rule supplies $D(T,E\vee F)=1$. Disjunction elimination then contradicts
(1). Therefore at least one of the two membership tests succeeds.
Finally, if $\bot\in T$, assumption and falsity elimination give
$D(T,C)=1$. This proves (5).
\end{proof}

The word ``maximal'' here refers only to the explicit finite condition (2).
There is no invocation of Zorn's lemma, an infinite enumeration of all
formulas, or a decision procedure for an unbounded proof search.

\subsection{Kripke semantics and the finite truth lemma}

A finite Kripke model consists of a nonempty finite partially ordered set
$(W,\leq)$ and an upward closed set $V(p)$ for every atom under consideration.
Atoms not listed in the finite input have the empty valuation. Define
forcing by
\[
\begin{array}{ll}
w\Forces p &\Longleftrightarrow w\in V(p),\\
w\not\Forces\bot,&\\
w\Forces B\wedge C &\Longleftrightarrow
                         (w\Forces B\ \hbox{and}\ w\Forces C),\\
w\Forces B\vee C &\Longleftrightarrow
                         (w\Forces B\ \hbox{or}\ w\Forces C),\\
w\Forces B\to C &\Longleftrightarrow
     \text{for every }v\geq w,\ v\Forces B\text{ implies }v\Forces C.
\end{array}
\]
For a finite coded frame these are Boolean computations. In the implication
clause, the universal quantifier is a loop over the finite list of worlds.
The formula recursion can be implemented by a bounded bottom-up scan of
its parse tree.

\begin{lemma}[Persistence and soundness]
\label{lem:ipc-soundness}
If $w\leq v$ and $w\Forces A$, then $v\Forces A$. Every $\IPC$ derivation
is valid in every finite Kripke model: if $\Gamma\vdash A$ and $w$ forces
all members of $\Gamma$, then $w\Forces A$.
\end{lemma}
\begin{proof}
For persistence, induction on formulas handles atoms by upward closure,
falsity trivially, and conjunction and disjunction componentwise. If $w$
forces $B\to C$ and $v\geq w$, every extension of $v$ is also an extension
of $w$, which proves the implication case.

For soundness, traverse the finite proof. Assumption, weakening, cut, falsity,
and the conjunction rules follow directly from the definitions. A disjunction
forced at $w$ gives one of its two forced disjuncts, so the appropriate
discharged branch proves the conclusion of disjunction elimination.
For implication introduction, let $v\geq w$ force the discharged assumption
$B$. Persistence carries every member of $\Gamma$ from $w$ to $v$, and the
premise derivation yields $v\Forces C$. This proves $w\Forces B\to C$.
Implication elimination uses the reflexive extension $w\leq w$.
These checks also cover proofs whose intermediate formulas do not belong
to the input subformula set $S$.
\end{proof}

Form the \emph{canonical finite frame} whose worlds are all subsets
$T\subseteq S$ satisfying the following three bounded conditions:
\begin{enumerate}
\item if $D(T,E)=1$, then $E\in T$, for every $E\in S$;
\item if $E\vee F\in T$, then $E\in T$ or $F\in T$;
\item $\bot\notin T$.
\end{enumerate}
Order these worlds by inclusion, and put $T\Forces p$ exactly when $p\in T$.
The frame is obtained by scanning all $2^m$ masks. When it is used to
refute a failed judgement, Lemma~\ref{lem:ipc-prime-extension} supplies at
least one world, so the required nonemptiness is established. We do not
need to assert nonemptiness before that hypothesis is available.

\begin{theorem}[Finite canonical truth lemma]
\label{thm:ipc-truth-lemma}
For every canonical world $T$ and every $E\in S$,
\[
                      T\Forces E\quad\Longleftrightarrow\quad E\in T.
\]
\end{theorem}
\begin{proof}
Use structural induction on $E$. The atomic case is the valuation definition;
the falsity case is condition (3). If $B\wedge C\in T$, assumption,
conjunction elimination, and condition (1) put both $B$ and $C$ in $T$.
Conversely, membership of both gives their conjunction by introduction and
condition (1). Apply the induction hypotheses. For disjunction, the forward
membership direction is condition (2), and the reverse direction follows
from disjunction introduction and condition (1).

Now let $E=B\to C$. If $E\in T$ and $U\supseteq T$ is a canonical world
forcing $B$, the induction hypothesis gives $B\in U$. Both $E$ and $B$
are assumptions in context $U$, so implication elimination puts $D(U,C)=1$.
Condition (1) and the induction hypothesis give $U\Forces C$. This proves
$T\Forces B\to C$.

Conversely, decide whether $B\to C\in T$. If it is absent, then
$D(T\cup\{B\},C)=0$: otherwise implication introduction would give
$D(T,B\to C)=1$, and condition (1) would put $B\to C$ in $T$.
Apply Lemma~\ref{lem:ipc-prime-extension} to $X=T\cup\{B\}$ and target $C$.
It produces a canonical world $U\supseteq T$ containing $B$ and avoiding
$C$. Assumption implies $C\notin U$. By the induction hypotheses,
$U\Forces B$ and $U\not\Forces C$, so $T\not\Forces B\to C$.
Thus whenever $T$ forces the implication, the membership test must have
succeeded. This proves the remaining direction constructively.
\end{proof}

\begin{theorem}[Proof or finite countermodel]
\label{thm:ipc-proof-or-model}
For every well-formed finite input $\Gamma_0\vdash A_0$, a primitive recursive
procedure returns exactly one of these outcomes, with correctness and the
stated bounds provable in $\ipra$:
\begin{enumerate}
\item an $\IPC$ derivation of $A_0$ from $\Gamma_0$; or
\item a finite Kripke model with at most $2^m$ worlds, and a distinguished
world forcing every member of $\Gamma_0$ but not $A_0$.
\end{enumerate}
\end{theorem}
\begin{proof}
Compute $D$. If $D(\Gamma_0,A_0)=1$, extract the proof from
Lemma~\ref{lem:ipc-finite-closure}. Otherwise extend $\Gamma_0$ to a
canonical world $T$ avoiding $A_0$. The truth lemma makes $T$ the required
counterexample, inside the frame of all canonical worlds. That frame has
at most $2^m$ worlds. Soundness shows that the two outcomes cannot both be
correct.

All loops are bounded: extraction of $S$ by the input parse trees, saturation
by $N$, greedy extension by $m$, and world enumeration by $2^m$. The
invariants used above can each be represented by a single Boolean condition
on finite arrays; bounded conjunctions are primitive recursive products of
bits. The inductions are on these arrays' construction stages, their bounded
counts, the parse-tree scan, and the proof scan. Thus they are instances of
quantifier-free induction after the primitive recursive characteristic
functions have been introduced. No induction on an unbounded semantic
assertion, reflection principle, or unbounded search is used.
\end{proof}

Consequently $\IPC$ derivability for finite contexts is decidable; it is
equivalent to validity in every finite Kripke model; and a failed sequent
has a countermodel within the displayed bound. In fact it suffices to test
all coded frames of at most $2^m$ worlds on the finitely many input atoms.
One such frame uses at most $q^2+aq$ Boolean entries for its order and atomic
valuation, where $q\leq2^m$ and $a\leq m$. Hence enumeration of all candidate
frames also has an explicit primitive recursive bound. The canonical
construction is preferable because it produces a witness directly.

There is also a useful syntactic consequence. Any ordinary $\IPC$ proof
of the input sequent, even if it uses arbitrary intermediate formulas,
can be replaced by a proof in the displayed finite restriction. Indeed,
compute the decision outcome. A countermodel would contradict soundness
of the supplied proof, so the proof outcome occurs. This proves a bounded
subformula result \emph{with cut allowed}; it is not an assertion that the
algorithm has performed cut elimination.

\subsection{Finite Heyting algebras}

For a finite poset $W$, let $\mathcal U(W)$ be its upward closed subsets.
This family is explicitly enumerable among the $2^{|W|}$ masks. Give it
the operations
\[
0=\varnothing,\qquad 1=W,\qquad
U\mathbin{\wedge}V=U\cap V,\qquad
U\mathbin{\vee}V=U\cup V,
\]
\[
U\Rightarrow V=
\{w\in W:\text{ every }v\geq w\text{ that belongs to }U
                              \text{ also belongs to }V\}.
\]

\begin{theorem}[Finite Heyting semantics and completeness]
\label{thm:ipc-heyting}
These operations make $\mathcal U(W)$ a finite Heyting algebra. The value
of a formula in this algebra is exactly its set of forcing worlds. Every
$\IPC$ theorem is valid in every coded finite Heyting algebra. Conversely,
if a finite sequent with $m$ relevant subformulas is not derivable, it fails
in a finite Heyting algebra with at most $2^{2^m}$ elements. All constructions
and implications are formalizable in $\ipra$.
\end{theorem}
\begin{proof}
Union and intersection preserve upward closure and satisfy the bounded
distributive lattice identities by membership tests. The set $U\Rightarrow V$
is upward closed because an extension of an extension is an extension.
For upward closed $X$, prove the defining adjunction
\[
             X\cap U\subseteq V
             \quad\Longleftrightarrow\quad
             X\subseteq U\Rightarrow V.
\]
In the forward direction, if $x\in X$ and $y\geq x$ lies in $U$, upward
closure puts $y\in X$, and the assumed inclusion puts $y\in V$.
For the reverse direction, apply the implication condition at $y=x$.
Thus this is a Heyting algebra. Formula induction identifies algebraic
values with forcing sets, including implication by its definition.

For an arbitrary finite Heyting algebra, write $g$ for the meet of the
values of the assumptions, with empty meet $1$. Each proof rule preserves
$g\leq[A]$. For example, implication introduction is exactly the adjunction
$g\wedge a\leq b$ iff $g\leq a\Rightarrow b$.
Implication elimination uses $a\wedge(a\Rightarrow b)\leq b$.
For disjunction elimination, $g\leq a\vee b$ and the two branch inequalities
$g\wedge a\leq c$, $g\wedge b\leq c$ imply
\[
g=g\wedge(a\vee b)=(g\wedge a)\vee(g\wedge b)\leq c.
\]
Cut uses $g\leq a$ and $g\wedge a\leq b$; the other rules follow directly
from the lattice laws. This proves soundness by a finite proof scan.

For completeness, use the countermodel from
Theorem~\ref{thm:ipc-proof-or-model}. Its distinguished world belongs to
the meet of the assumption values but not to the conclusion value in
$\mathcal U(W)$. Thus the required inequality fails. The algebra contains
at most $2^{|W|}\leq2^{2^m}$ elements. All algebra axioms and operations
in this argument are finite tables and bounded checks.
\end{proof}

\subsection{Disjunction extraction, including Harrop assumptions}

A \emph{Harrop formula} is generated by
\[
H::=p_i\mid\bot\mid H\wedge H\mid B\to H,
\]
where $B$ is arbitrary. In particular every negation is Harrop. We prove
a primitive recursive extraction result. Stronger polynomial bounds are
known \cite{buss-pudlak2001}; no such bound is needed or claimed here.

\begin{lemma}[Gluing two rooted models]
\label{lem:ipc-gluing}
Let $M_1,M_2$ be disjoint finite rooted Kripke models, with roots $r_1,r_2$.
Adjoin a fresh root $r$ below every world of both models. Give an atom to
$r$ exactly when both $r_1$ and $r_2$ force it. Forcing at the old worlds
is unchanged. If a Harrop formula is forced at both $r_1$ and $r_2$, then
it is forced at $r$.
\end{lemma}
\begin{proof}
The new atomic valuation is persistent: an atom at $r$ belongs to both old
roots and hence to all their extensions. An old world has exactly its old
extensions, so induction on formulas shows that its forcing relation is
unchanged.

For the final assertion, induct on the displayed Harrop grammar. Atoms
hold by the definition at $r$; the falsity case has an impossible premise;
conjunction follows componentwise. Suppose $H=B\to K$ is forced at both
old roots. Every old extension forcing $B$ forces $K$ already. It remains
to consider $r$ itself: if $r\Forces B$, persistence gives $r_i\Forces B$
for both $i$, hence $r_i\Forces K$. The induction hypothesis for the
strictly smaller Harrop formula $K$ gives $r\Forces K$. This is precisely
the implication condition at the new root.
\end{proof}

\begin{theorem}[Constructive disjunction property]
\label{thm:ipc-disjunction}
If all formulas in the finite context $\Gamma$ are Harrop, an $\IPC$ proof
of $\Gamma\vdash A\vee B$ can be transformed by a primitive recursive
procedure into either a proof of $\Gamma\vdash A$ or a proof of
$\Gamma\vdash B$, together with the tag identifying the chosen disjunct.
In particular, $\IPC$ has the disjunction property for theorems.
\end{theorem}
\begin{proof}
Run the proof-or-model procedure for both $\Gamma\vdash A$ and
$\Gamma\vdash B$. If either returns a proof, output it and its tag.
Otherwise obtain two countermodels with distinguished worlds. Restrict
each model to the upward cone of its distinguished world; this preserves
forcing there and makes that world a root. Glue the two models as in
Lemma~\ref{lem:ipc-gluing}. Every assumption holds at the fresh root because
it is Harrop and holds at both old roots.

The new root cannot force $A$: persistence would carry $A$ to the old root
of its countermodel. The same argument excludes $B$. Thus the fresh root
does not force $A\vee B$, contradicting soundness of the supplied proof.
Consequently one of the two decidable proof searches succeeds. An
implementation may return a default value on the impossible third branch;
the theorem's proof-checking hypothesis proves that this branch is never
reached on valid input. All searches, cone restrictions, and gluing
operations are bounded and primitive recursive.
\end{proof}

\paragraph{Disjunction cannot be defined by conjunction, implication and falsity.}
More generally no Harrop formula $H$ is intuitionistically equivalent to
$p\vee q$ for distinct atoms $p,q$. To see this, take two one-world models,
one with exactly $p$ true and one with exactly $q$ true. If both directions
of the proposed equivalence had proofs, soundness would make $H$ true at
both worlds, since $p\vee q$ is true there. Glue these models as above.
The Harrop gluing lemma makes $H$ true at the new root, while both $p$ and
$q$ are false there. Soundness of $H\to(p\vee q)$ gives a contradiction.
This is uniform in the finite code of $H$ and its alleged equivalence
proof. Every formula built from atoms, $\bot,\wedge,\to$ is Harrop, so the
claimed nondefinability follows. This is a failure of intuitionistic
definability; the classical definition in Theorem~\ref{thm:three-axiom}
remains valid.

The deduction rule gives the familiar admissible Harrop rule as a further
explicit consequence:
\[
\frac{\vdash\neg A\to(B\vee C)}
     {\vdash(\neg A\to B)\vee(\neg A\to C)}.
\]
Indeed, apply disjunction extraction under the Harrop assumption $\neg A$,
discharge that assumption in the selected proof, and introduce the outer
disjunction. Admissibility means a transformation of theorem proofs. It
does not assert that the implication from the premise formula to the
conclusion formula is itself an $\IPC$ theorem.

\subsection{Small countermodels and exact failures of classical laws}

The following finite models certify nonderivability by
Lemma~\ref{lem:ipc-soundness}. They also show why a classical truth table
cannot serve as an intuitionistic completeness test.

\paragraph{A two-world chain.}
Take $r<t$ and let $p$ hold only at $t$. At $r$, neither $p$ nor $\neg p$
holds: $p$ is absent there, while $t$ witnesses the failure of $\neg p$.
No world forces $\neg p$, so $r$ does force $\neg\neg p$. Therefore both
\[
              p\vee\neg p,
              \qquad \neg\neg p\to p
\]
fail at $r$. The implication $\neg p\to p$ holds at $r$ because no extension
forces $\neg p$; hence the Peirce instance
\[
                      ((p\to\bot)\to p)\to p
\]
also fails there. Finally $p\to p$ holds everywhere but $\neg p\vee p$
fails at $r$, refuting the general replacement of $A\to B$ by
$\neg A\vee B$.

\paragraph{A three-world fork.}
Take $r<u$, $r<v$, with $u,v$ incomparable and terminal. Let $p$ hold only
at $u$ and $q$ only at $v$. The following formulas fail at $r$:
\[
\begin{gathered}
\neg p\vee\neg\neg p,\qquad
(p\to q)\vee(q\to p),\\
\neg(p\wedge q)\to(\neg p\vee\neg q),\qquad
\neg\neg(p\vee q)\to(\neg\neg p\vee\neg\neg q).
\end{gathered}
\]
For the first, $u$ refutes $\neg p$ at $r$, and $v\Forces\neg p$ refutes
$\neg\neg p$ there. For the second, $u$ and $v$ refute the respective
implications. For the third, no world forces $p\wedge q$, while neither
negated atom holds at $r$. For the fourth, every terminal world forces
$p\vee q$, so no world forces its negation; nevertheless each individual
double negation fails at $r$ because of the opposite branch.

On the same fork, assign an additional atom $a$ to both terminal worlds
and to no other world. Then
\[
                (a\to(p\vee q))\to
                   ((a\to p)\vee(a\to q))
\]
fails: the antecedent holds at $r$, while each disjunct of the consequent
has a different terminal counterexample.

\paragraph{An admissible rule need not be an axiom.}
For the Harrop rule above, use a root with three terminal successors, with
exactly $p$, exactly $q$, and exactly $s$ true respectively. At the root,
$\neg p\to(q\vee s)$ holds: the two terminal worlds forcing $\neg p$
force the required disjunction, and the remaining worlds do not force
$\neg p$. Neither $\neg p\to q$ nor $\neg p\to s$ holds at the root.
Thus
\[
(\neg p\to(q\vee s))\to
       ((\neg p\to q)\vee(\neg p\to s))
\]
is not an $\IPC$ theorem even though the associated rule is admissible.

These are uniform arithmetical nonderivability statements about the
specified propositional proof checker, proved by concrete finite
evaluations and soundness. They are not claims that $\ipra$ refutes every
arithmetical substitution instance of the displayed formulas.

\subsection{Lean implementation obligations for this construction}

The implementation should retain two distinct levels. An inductive
\texttt{Formula} and an inductive finite derivation type provide readable
Lean specifications. Natural-number codes, Boolean tables, and fixed
iteration bounds provide the primitive recursive objects interpreted in
$\ipra$. A Lean proof about the readable specification becomes an
$\ipra$ result only after the coding and quantifier-free invariant bridge
has been proved.

The necessary definitions are: formula coding and decoding; subformula
enumeration with duplicate removal; context masks and their union,
inclusion and insertion operations; an index for each of the $N$ sequents;
the twelve-rule local checker; the inflationary table update; the
$N$-fold iteration; first-appearance records and their proof checker; the
greedy extension; the predicate on masks defining canonical worlds;
finite frame validity; and the bottom-up Boolean forcing evaluator.

The main proof dependency order is:
\begin{enumerate}
\item Prove formula coding and subformula closure, and the cardinality
      $N=m2^m$ of the sequent table.
\item Prove the update is inflationary and that equality of two successive
      tables implies permanence. Prove stabilization at $N$ by the set-bit
      count, and prove soundness of first-appearance certificates.
\item Prove the five greedy-extension invariants in
      Lemma~\ref{lem:ipc-prime-extension}. These require only the closed
      table's weakening, cut, disjunction elimination, and falsity rules.
\item Prove persistence and proof soundness for arbitrary finite models.
      Prove the canonical truth lemma by induction on a formula's coded
      parse tree; use the greedy extension only in the implication case.
\item Combine the previous results into a total tagged proof-or-countermodel
      function, with certified output checkers and the explicit bounds.
\item Define upward-set algebra tables and prove the adjunction by finite
      membership reasoning. Define cone restriction and gluing; prove
      preservation and the Harrop induction, then the extraction theorem.
\end{enumerate}
No opaque axiom asserting completeness or decidability belongs in this
dependency graph. A small executable can test the supplied countermodels,
but such tests do not discharge any of the general Lean obligations.
